%*******************************************************
% Annexe C — Hypothèse d'invariance structurelle fonctionnelle
% Avril 2026 — Note méthodologique
%*******************************************************

\chapter{L'hypothèse d'invariance structurelle fonctionnelle~: trois formulations}\label{ch:annexe-invariance}

\noindent\emph{Cette note explicite l'hypothèse centrale qui sous-tend l'architecture de la thèse. Les chapitres~1 et~2 ne la formulent pas directement~--- le chapitre~1 construit le matériau historique par induction, le chapitre~2 en extrait les mécanismes~--- mais elle organise le passage de l'un à l'autre. Trois formulations de force croissante sont proposées. Chacune est assortie de ses conditions de validité et de son critère de réfutation.}

\bigskip

\section{La question}

Le chapitre~1 observe que des technologies de l'information séparées par un siècle et demi~--- le télégraphe de Morse, les calculatrices mécaniques, le système mécanographique d'IBM~--- produisent des effets économiques qui se laissent décrire dans un vocabulaire commun~: substitution conditionnée par la demande, irréversibilité croissante avec l'intégration systémique, détermination de la structure de marché par la forme du capital. Le chapitre~2 systématise ces régularités en dix dimensions et identifie des mécanismes causaux récurrents.

La question est~: cette récurrence est-elle un artefact de la grille d'analyse (on retrouve ce qu'on a mis), une régularité empirique sans explication causale, ou le signe d'une contrainte structurelle que les technologies de l'information partagent en raison de ce qu'elles \emph{font}~--- stocker, transmettre, traiter~?

\section{Version faible~: la régularité empirique}

\paragraph{Énoncé.} Les technologies de l'information, à travers leurs incarnations successives, produisent des mécanismes économiques récurrents. Les mêmes structures causales~--- substitution entre capital informationnel et capital physique, irréversibilité conditionnée par l'intégration systémique, détermination de la structure de marché par le mode d'appropriation du capital~--- s'observent du télégraphe au data center.

\paragraph{Statut.} C'est une proposition descriptive~: elle constate un pattern sans l'expliquer. Elle est du même ordre que les faits stylisés de \textcite{kaldor_capital_1961}~--- des régularités empiriques suffisamment robustes pour contraindre les modèles théoriques, mais qui n'impliquent pas de mécanisme générateur unique.

\paragraph{Conditions de validité.} La régularité doit être documentée dans au moins trois cas historiquement indépendants (c'est-à-dire non liés par une filiation technologique directe). Le chapitre~1 en fournit quatre~: le télégraphe, les calculatrices mécaniques, la mécanographie IBM, et le test contemporain du chapitre~2.

\paragraph{Critère de réfutation.} L'hypothèse serait affaiblie si l'on identifiait une technologie de l'information majeure dont les effets économiques ne se laissent pas décrire par les mécanismes du chapitre~2~--- c'est-à-dire une technologie qui produit une substitution, une irréversibilité ou une structure de marché qui ne relèvent d'aucune des dix dimensions. Le test le plus direct serait un cas où la forme du capital ne prédit pas le type d'effet.

\paragraph{Ce que cette version ne dit pas.} Elle ne dit pas \emph{pourquoi} les mécanismes récurrent. La récurrence pourrait tenir à une contrainte profonde, à une coïncidence historique, ou à un biais de sélection de l'observateur (les cas choisis dans le chapitre~1 partagent des propriétés parce qu'ils ont été choisis pour cela). La version faible est compatible avec les trois explications.

\section{Version intermédiaire~: la contrainte fonctionnelle}

\paragraph{Énoncé.} Les mécanismes économiques des technologies de l'information récurrent parce qu'ils sont contraints par les fonctions élémentaires que ces technologies remplissent~--- stocker, transmettre, traiter. Toute technologie de l'information doit réaliser au moins l'une de ces trois opérations pour produire un effet économique~; la manière dont elle le fait~--- le substrat matériel, l'architecture technique, le contexte institutionnel~--- change d'une époque à l'autre, mais la structure fonctionnelle persiste. Cette structure fonctionnelle délimite l'espace des effets économiques possibles.

\paragraph{Statut.} C'est une proposition explicative~: elle identifie la décomposition stocker/transmettre/traiter (S/T/P) comme le mécanisme générateur de la récurrence. Le chapitre~1 construit cette décomposition par induction (section~\ref{subsec:jacquard} pour le stockage sur carte perforée, section~\ref{sec:telegraph} pour la transmission, sections~\ref{sec:calculating} et~\ref{sec:mechanography} pour le traitement). Le chapitre~2 l'utilise implicitement lorsqu'il montre que la nature de la demande (coordination = transmission, traitement interne = processing, mémoire = stockage) conditionne le type d'effet.

\paragraph{Conditions de validité.} (i)~La décomposition S/T/P doit être exhaustive~: toute opération informationnelle doit pouvoir se décomposer en combinaison de ces trois fonctions. (ii)~La décomposition doit être discriminante~: deux technologies qui remplissent des fonctions différentes dans le triangle S/T/P doivent produire des effets économiques différents. (iii)~La décomposition doit être indépendante du substrat~: elle doit s'appliquer que le substrat soit mécanique (Jacquard), électrique (télégraphe), électronique (GPU) ou autre.

\paragraph{Critère de réfutation.} L'hypothèse serait réfutée si l'on identifiait~: (a)~une opération informationnelle qui ne se décompose pas en S/T/P~; ou (b)~deux technologies qui occupent la même position dans le triangle S/T/P, dans le même contexte de demande, mais qui produisent des effets économiques structurellement différents~; ou (c)~une technologie qui change de substrat matériel sans changer de position fonctionnelle et dont les effets économiques changent néanmoins de structure (ce qui montrerait que le substrat, et non la fonction, détermine l'effet).

\paragraph{Ce que cette version ne dit pas.} Elle ne dit pas que la contrainte fonctionnelle \emph{détermine} les effets économiques~--- elle dit qu'elle les \emph{contraint}. La demande, les institutions, la géographie et le contexte politique restent des variables qui modulent le résultat dans l'espace que la contrainte fonctionnelle délimite. La calculatrice et la tabulatrice remplissent la même fonction (traiter) mais produisent des effets différents parce que le mode d'appropriation diffère~--- et le mode d'appropriation est un choix institutionnel, pas une nécessité fonctionnelle. La contrainte fonctionnelle est une condition nécessaire, pas suffisante.

\section{Version forte~: la détermination par les propriétés physiques de l'information}

\paragraph{Énoncé.} Les propriétés physiques de l'information elle-même~--- la rivalité conditionnelle du signal (un canal occupé ne peut pas transmettre deux messages simultanément), la non-dégradation de la copie (l'information peut être reproduite sans perte), la dépendance au substrat matériel pour les trois fonctions (stocker requiert un support, transmettre requiert un médium, traiter requiert de l'énergie)~--- déterminent les formes économiques possibles des technologies de l'information. Les mécanismes identifiés dans les chapitres~1 et~2 ne sont pas seulement contraints par les fonctions~; ils sont \emph{nécessités} par les propriétés physiques du bien informationnel.

\paragraph{Statut.} C'est un argument de nécessité. Il est le plus séduisant et le plus dangereux. Il est séduisant parce qu'il fonde la récurrence dans la physique plutôt que dans l'histoire~: les mécanismes récurrent parce qu'ils \emph{doivent} récurrer. Il est dangereux parce qu'il flirte avec le déterminisme technologique~--- la position selon laquelle les propriétés de la technologie déterminent les formes sociales, que la tradition STS (Science and Technology Studies) a critiquée depuis \textcite{hughes_networks_1983} et que les travaux de Geels sur les transitions sociotechniques ont rendue difficilement défendable dans sa forme forte.

\paragraph{Conditions de validité.} La version forte n'est défendable que si l'on peut montrer que les propriétés physiques de l'information contraignent les formes économiques \emph{indépendamment} du contexte institutionnel et géographique. Cela exigerait de documenter au moins un cas où le même mécanisme opère dans des contextes institutionnels radicalement différents~--- par exemple, un marché capitaliste et une économie planifiée~--- avec le même résultat structurel.

\paragraph{Critère de réfutation.} L'hypothèse serait réfutée si l'on identifiait deux contextes institutionnels différents dans lesquels la même technologie, remplissant la même fonction, avec les mêmes propriétés physiques, produisait des effets économiques structurellement différents. Le cas de la Chine contemporaine (régulation étatique de l'IA, marché intérieur fermé, infrastructure de calcul sous embargo) pourrait être un terrain de test~: si les mécanismes du chapitre~2 ne s'y appliquent pas, la version forte est en difficulté.

\paragraph{Ce que cette version implique.} Si elle était vérifiée, elle impliquerait que les modèles économiques qui traitent l'ICT comme un facteur de production homogène ne sont pas seulement empiriquement insuffisants (ce que les versions faible et intermédiaire montrent)~--- ils sont \emph{théoriquement mal fondés}, parce qu'ils ignorent des contraintes physiques qui déterminent la forme des effets. C'est une thèse forte. Elle n'est pas nécessaire au fonctionnement de la présente thèse, qui repose sur la version intermédiaire.

\section{Position de la thèse}

La présente thèse adopte la version intermédiaire comme hypothèse de travail. Le chapitre~1 fournit le matériau inductif~; le chapitre~2 teste la portée des mécanismes par confrontation systématique avec dix dimensions économiques~; le chapitre~3 examine les conditions sous lesquelles ces mécanismes changent de régime (le couplage énergie-information comme rupture dans le gradient d'irréversibilité). La version faible est un résultat acquis dès le chapitre~1. La version forte est une conjecture qui dépasse le périmètre de cette thèse mais qui en constitue un prolongement naturel~--- en particulier si les cas géographiquement diversifiés (Afrique subsaharienne, Chine, Europe) du test contemporain du chapitre~2 confirment la persistance des mécanismes à travers les contextes institutionnels.
